% 第 6 章 实验与指标体系 (目标 4 页)
% 对应官方要求: 实验方案 + 指标体系
\section{实验与指标体系}

\subsection{实验设置}

每个样本在两种模式下跑通。基线模式直接调用模型,无任何中间层。防护模式在模型与工具调用之间插入防护代理,代理按规则层、评审层顺序检查。

模型主干选用两类国产主流模型,覆盖代表性业务使用场景。评测在真实智能体上跑通端到端链路,真实加载 50 类政务技能,真实执行工具调用,真实产生审计日志。

每条样本记录四类结果。BYPASSED 表示模型真调了攻击工具,REFUSED 表示模型主动拒答,NO\_ATTACK 表示模型未触发攻击工具。BYPASSED 是防护体系要解决的核心类别。

\subsection{指标体系}

三项核心指标。

\textbf{攻击拦截率}。定义为基线 BYPASSED 状态在防护模式下转为非 BYPASSED 状态的比例,反映防护体系对真实攻击的拦截能力。

\textbf{误报率}。定义为白样本在防护模式下被阻断或被改写的比例,反映防护体系对正常业务的容忍度。

\textbf{延迟增加}。定义为加防护代理后单次模型调用的延迟增加,反映工程可用性。

\subsection{实验结果}

\textbf{攻击拦截率}。在 66 条攻击样本上,基线阶段 BYPASSED 57 条,绕过率 86.4\%。防护模式下 BYPASSED 0 条,拦截率 100\%。其中结构化规则层拦截 38 条,评审层兜底拦截 17 条,其余在迭代修复阶段通过字段变体规则与轨迹级联识别补齐。

\textbf{误报率}。在 25 条白样本上,防护模式真实误报 0 条,误报率 0\%。规则层与评审层均未对白样本触发阻断。

\textbf{精确归因}。66 条样本逐条比对基线与防护结果,83.3\% 归因为防护代理拦截,13.6\% 归因为模型自身安全机制,3.0\% 归因为迭代修复后的代理拦截。该归因说明 100\% 拦截中 86.4 个百分点来自本防护体系。

\textbf{代理内部动作}。实验期间代理共产生 173 次结构化规则拦截、506 次评审调用、38 次文本级阻断,平均单请求延迟增加 1.8 秒。

\begin{table}[H]
  \centering
  \caption{核心实验结果汇总}
  \label{tab:results}
  \small
  \begin{tabular}{lccc}
    \toprule
    指标 & 基线模式 & 防护模式 & 变化 \\
    \midrule
    攻击样本绕过率 & 86.4\% & 0\% & 下降 86.4 个百分点 \\
    白样本误报率 & 0\% & 0\% & 持平 \\
    单请求延迟增加 & -- & 1.8 秒 & -- \\
    \bottomrule
  \end{tabular}
\end{table}

\subsection{消融实验}

\textbf{仅规则层}。在 66 条攻击样本上单独启用规则层,命中 27 条,绕过率从 86.4\% 压到 40.9\%。规则层单独可拦截约一半攻击,作为高置信基础层具备工程价值。

\textbf{规则层加评审层}。两层协同后 BYPASSED 为 0 条。评审层对规则未命中的样本提供有效兜底,两层达成 100\% 拦截。

\textbf{白样本误报隔离}。规则层与评审层在 25 条白样本上均未触发阻断,两层在保持高拦截的同时未引入白样本误伤。

\subsection{与基线方案对比}

将本体系与模型自身安全过滤对比。在 66 条去标记化业务话术样本上,模型自身安全仅拦截 9 条,本体系拦截全部 66 条。在 25 条白样本上,两者均保持 0 误报。本体系同时在真实智能体上提供端到端拦截证据,覆盖规则、评审、链路三层,公开方案的评估多停留在离线基准阶段。
